\section{Scheduler Benchmark Methodology}
\label{sec:methodology}

\subsection{Scheduler Panel}
\label{sec:panel}

The PRIMARY panel comprises 11 policies, executed alongside 2 additional
\texttt{STYLE\_APPROXIMATION} policies used only for a mandatory robustness
check (13 policies executed in total): a classical arrival-order baseline
(\texttt{fifo}, also the load-calibration reference policy), a
deadline-aware policy (\texttt{edf}), a slack/urgency-aware policy
(\texttt{least\_laxity\_first}), a service-length-aware policy
(\texttt{estimated\_service\_time\_first}), a fairness-aware policy
(\texttt{weighted\_fair\_share}), a KV/memory-pressure-aware policy
(\texttt{kv\_constrained\_online}), an admission-control/SLO-guard policy
(\texttt{admission\_control}), and faithful reimplementations of four
external mechanisms: vLLM-style token-budget continuous batching in both
its pre-chunked-prefill and chunked-prefill forms
(\texttt{vllm\_faithful}, \texttt{vllm\_chunked\_prefill\_faithful}),
Sarathi-style chunked prefill (\texttt{sarathi\_faithful}), and an
SLO-aware scheduling policy in the SLAI/RAD family
(\texttt{slai\_faithful}). The two \texttt{STYLE\_APPROXIMATION} robustness
policies are a token-budget-proxy style approximation
(\texttt{vllm\_style\_token\_budget}) and an SLO-guard-proxy style
approximation (\texttt{scorpio\_style\_slo\_guard}); two further faithful
reimplementations requiring disaggregated/cross-instance simulator
semantics (\texttt{distserve\_faithful}, \texttt{llumnix\_faithful}) are
executed only in a secondary stratum, never pooled with the single-GPU
colocated PRIMARY leaderboard, and one candidate
(\texttt{apt\_serve\_faithful}) is excluded entirely as
\texttt{scaffolding\_only} -- not a validated reimplementation at the time
of this freeze.

A policy enters the PRIMARY panel only if, at the time it was added to the
comparability audit (frozen the day before Stage 0 executed), it satisfies
an outcome-independent selection rule: its mechanism predates this
protocol's freeze, it occupies a mechanism family not already saturated by
an existing panel member, it has adequate implementation fidelity and
passes its own unit tests, it is simulator-semantically compatible with the
single-GPU colocated execution model, and it is independently reproducible
from a cited reference. No policy was added to or removed from the panel in
response to any Stage-0 outcome (\S\ref{sec:stage0}); the panel document
predates Stage-0 execution by a full day, and mechanism families identified
as under-implemented (preemptive/multilevel scheduling in the FastServe
sense, JITServe, LMetric) were deliberately left as documented future-work
candidates rather than fast-tracked into the panel after seeing Stage-0's
BurstGPT collapse.

\subsection{Fidelity Taxonomy}
\label{sec:fidelity}

Every executed policy is labeled with one of five fidelity classes:
\texttt{FAITHFUL\_EXTERNAL} (a validated reimplementation of a specific
published mechanism), \texttt{REPOSITORY\_NATIVE\_CLASSICAL} (a standard
scheduling discipline implemented natively in this simulator, not claimed
as a reproduction of any specific paper), \texttt{SIMULATOR\_PROXY} (a
mechanism-motivated policy without a single canonical external paper),
\texttt{STYLE\_APPROXIMATION} (an approximate, style-level proxy for a
published mechanism, explicitly excluded from the PRIMARY headline panel
and used only for robustness checks), and \texttt{scaffolding\_only} (an
unaudited, not-yet-validated implementation, excluded entirely). We do not
claim that any \texttt{STYLE\_APPROXIMATION} or \texttt{SIMULATOR\_PROXY}
policy is an official or validated reimplementation of its namesake
published system.

\subsection{Execution Model}
\label{sec:execution-model}

Each (window, policy, operating-region) cell is executed as a fully
deterministic discrete-event simulation: given identical inputs, the
simulator produces bit-identical outputs, with no consumed randomness
(verified by a repeated-run determinism test). This determinism is what
permits the paired design of \S\ref{sec:unit}: two policies' outcomes on
the same window and region differ only because of the policy's own
scheduling decisions, not because of independent sampling noise.

\subsection{Stage-0 Discriminability Pilot (Historical)}
\label{sec:stage0}

Before the Pilot-V2 protocol described in the remainder of this section was
frozen, a preregistered discriminability pilot (``Stage 0'') executed a
real, 1,080-cell matrix (array 1213964; three sources, three load regions, a
six-policy PRIMARY subset) to check, before committing to the full
protocol, that the panel would yield non-trivial scheduler differentiation
under a fixed load protocol without post-hoc cherry-picking. \textbf{This
is historical pilot evidence, not the confirmatory Pilot-V2 result}, and it
never measured cross-source ranking portability -- that question was and
remains explicitly out of Stage-0's scope. Stage 0's own five-criterion
discriminability analyzer returned four passing criteria (77.8\% of
(source, window, load-region) triples non-tied, 86.7\% of windows
non-tied, no universal policy collapse, and 100\% of non-tied cells showing
meaningful metric variation) and one genuine failure: BurstGPT contributed
only 14.3\% of non-tied conditions, against a 15\% floor, versus
$\sim$42.9\% each for the other two sources -- a result reported as
\texttt{STAGE0\_NO\_GO}, robust to the completion-conditioned-metric
zero-completion convention used. A follow-up diagnostic traced this to a
mechanism-level explanation: differentiation was strongly source-dependent,
attributable to a near-total collapse of five of the six Stage-0 policies
into pairwise bit-identical outcomes on BurstGPT's structurally
short-prompt, near-constant-output traffic -- not to BurstGPT receiving
weaker or fewer calibrated load conditions than the other sources. This
motivated the Pilot-V2 protocol's larger, mechanism-diverse 13-policy panel
and denser six-region load grid (\S\ref{sec:calibration}), applied
symmetrically to every source, never to BurstGPT alone. Detailed Stage-0
mechanism analysis is deferred to an appendix; it is not restated here as a
ranking-portability finding.

\subsection{Operating-Region Calibration}
\label{sec:calibration}

A FIFO-only calibration pass establishes six operating regions --
\texttt{LOW} (load factor 0.5), \texttt{PRE\_KNEE} (0.8), \texttt{KNEE}
(1.0), \texttt{POST\_KNEE} (1.1), \texttt{OVERLOAD} (1.2), and
\texttt{HIGH\_PRESSURE} (1.5) -- across all 120 frozen windows (720
calibration cells total; 720 valid, deterministic, 0 failed/missing/
duplicated). We report this as calibration provenance, not a comparative
scheduler result: FIFO is the calibration reference policy, and the
calibration's role is to fix load-factor-to-region mappings before any
other policy is executed against them, not to characterize FIFO's relative
performance. As a calibration characteristic worth preserving, 301 of the
720 calibration-region records (93 BurstGPT, 101 Bailian/Qwen, 107 Azure)
recorded zero completions under FIFO at that region's load factor; this is
a property of the calibration design at those specific (window, region)
combinations, the frozen calibration validity gate passed regardless
(\texttt{PHASE11\_CALIBRATION\_VALID = YES}), completion-conditioned metrics
are simply undefined (not zero) in those records
(\S\ref{sec:metrics}), and no calibration redesign was performed after
observing this count.

\subsection{Metrics}
\label{sec:metrics}

Three metrics are always defined for every completed cell: completion
fraction, arrival-normalized weighted goodput (ANWG, the primary metric),
and weighted completion fraction. A second group of metrics is defined
only conditional on at least one request completing: SLO-violation rate,
weighted goodput, mean/p95/p99 latency, and request/token throughput are
each \texttt{NaN} if and only if completion fraction is exactly zero for
that cell. Mean/p95 time-to-first-token is subject to a stricter,
separately checked precondition (at least one completed request recorded a
first-token time), because TTFT can be undefined even when completion
fraction is positive. No undefined conditional-metric value is ever imputed
as 0 or 1, in either direction, and never after inspecting which policy or
source it would affect. Where a policy's value is undefined for a given
(window, region, metric) condition, that policy is excluded from the
ranking-derived statistic (Kendall's tau-b, Spearman's rho, top-$k$
overlap, pairwise reversal detection) for that specific condition and
metric only -- never scored as a loss or a tie, and never causing the whole
condition to be dropped for policies that did complete. Every
ranking-derived table reports its exclusion count
(\texttt{n\_conditions\_excluded\_for\_undefined\_metric}) alongside the
statistic itself, since a source-concentrated exclusion rate is itself a
reportable finding under RQ4, not a footnote.

\subsection{Paired Design and Repetitions}
\label{sec:paired}

Because the simulator is fully deterministic (\S\ref{sec:execution-model}),
a single execution per (window, policy, region) cell is sufficient;
ranking-uncertainty quantification (\S\ref{sec:results}) instead comes from
resampling over the independent workload windows themselves (block
bootstrap), not from repeated stochastic runs of the same cell.

\subsection{Mechanism Telemetry}
\label{sec:telemetry}

Every executed cell additionally records a seven-field mechanism-activation
telemetry summary: queue-depth (mean/max), batch-saturation (mean/max
fraction of \texttt{max\_active\_sequences} in use), prefill/decode
contention fraction, KV-occupancy (mean/max fraction of \texttt{max\_kv\_tokens}
in use), admission-control activation count, preemption/reorder event
count, and token-budget saturation fraction. Each field is always defined:
a value of zero means either that the corresponding mechanism is genuinely
absent for that policy (for example, preemption count is exactly zero for
every policy that never issues a preempt/swap/migrate action, which is most
of the panel) or that the underlying simulator mode does not implement a
separate prefill/decode phase or token-budget model -- not an
instrumentation failure. This telemetry is collected for every cell but is
not used to support any mechanism-level finding in this manuscript; RQ4's
telemetry-to-reversal association analysis (\S\ref{sec:results}) is
pre-specified and will only be run once Pilot-V2 outcomes exist.

\subsection{Robustness Analyses}
\label{sec:robustness}

Every headline ranking-portability result (\S\ref{sec:results}) is
pre-specified to be checked against a fixed set of robustness analyses: the
11-policy high-fidelity subset (PRIMARY panel excluding both
\texttt{STYLE\_APPROXIMATION} policies), leave-one-source-out re-estimation,
window-size sensitivity, an alternate completion-conditioned-metric
convention, the full six-region load grid, a temporal-split sensitivity
check, and mechanism-family-exclusion checks. These robustness analyses are
part of the frozen protocol; their results are reported alongside
\S\ref{sec:results} once Pilot-V2 executes (\S\ref{tab:robustness}).
